\section{An original native affine chart is the reconstructed whole-field model chart}
\label{sec:projective-native-model-chart-algebra}

\begin{theorem}[Literal original affine sections equal the original coordinate algebra]
\label{thm:projective-native-chart-section-algebra}
\lean{EnvelopeProperty.ProjectiveNativeModelChartAlgebra.sectionSubalgebra_eq_adjoin_coordinates}
\leanok
\uses{thm:projective-chart-coefficient-algebra}
\source{bj22-trivval:sec:tc}
Let $f:X\hookrightarrow\mathbf P^m_k$ be an actual original closed
immersion of an integral scheme and let $U=f^{-1}D_+(X_i)$ be nonempty.
The literal generic-germ image of its entire original section algebra
$\Gamma(X,U)$ in the same native $K(X)$ is exactly
$k[X_j/X_i\mid 0\leq j\leq m]$, using the original pulled coordinate
sections. The original coefficient action is the structure map
$f$ followed by the original projective coefficient projection.
\end{theorem}
\begin{proof}
Every original pulled coordinate ratio is the generic germ of an actual
section on $U$, so its generated subalgebra is contained in the literal
native section image. Conversely the actual closed immersion gives
surjectivity from the whole original projective chart section ring onto
$\Gamma(X,U)$. The original polynomial-chart algebra isomorphism
identifies that source with the polynomial ring on the nonpivot ratios.
Its generic pullback is a genuine homomorphism over the original $k$,
and it sends every original polynomial variable to its literal pulled
generic coordinate ratio. Polynomial induction therefore puts the image
of every original section in the coordinate subalgebra. These two
containments give equality as subalgebras of the same original field.
\end{proof}

\begin{theorem}[Whole native affine chart and the actual reconstructed model chart]
\label{thm:projective-native-model-chart-algebra}
\lean{EnvelopeProperty.ProjectiveNativeModelChartAlgebra.modelChartAlg_eq_adjoin_coordinates, EnvelopeProperty.ProjectiveNativeModelChartAlgebra.nativeSectionsEquivModelChart, EnvelopeProperty.ProjectiveNativeModelChartAlgebra.nativeSectionsEquivModelChart_apply, EnvelopeProperty.ProjectiveNativeModelChartAlgebra.nativeChartIso, EnvelopeProperty.ProjectiveNativeModelChartAlgebra.nativeChartIso_structure}
\leanok
\uses{thm:projective-native-chart-section-algebra,thm:projective-chart-whole-native-field-model}
\source{bj22-trivval:sec:tc}
For the actual rank-one homogeneous model $P$ reconstructed over the
entire same native $K(X)$ from this embedding and these literal ratios,
the model chart with original pivot $i$ has exactly the same coordinate
subalgebra. Thus the entire original $\Gamma(X,U)$ is isomorphic over
the same $k$ to this actual model chart algebra, with every original
section mapped to its literal original generic germ. The original open
scheme $U$ is actually isomorphic to the spectrum of that reconstructed
model chart algebra, retaining the literal original coefficient
triangle over $k$. No scheme, section-ring or coordinate comparison
is supplied as a premise.
\end{theorem}
\begin{proof}
The original pivot coefficient in the reconstructed model is one. Its
actual chart dehomogenization therefore sends each homogeneous coordinate
$f_jT$ to the same literal original ratio $f_j$. The model homogeneous
algebra is generated by these original coordinates, so its chart algebra
is exactly their original generated subalgebra. Combining this equality
with the preceding whole native section-image equality and the actual
injectivity of generic germs on an integral scheme gives the claimed
whole original algebra isomorphism. Its underlying field value is
literally the original generic germ, by construction.

The original projective coordinate open is affine and its inverse image
under the actual closed immersion is affine. Its canonical native
isomorphism to the spectrum of its original entire section ring,
followed by the spectrum of the inverse proved ring isomorphism, gives
the actual whole scheme isomorphism to the reconstructed model chart.
The proved algebra equivalence preserves the original $k$ action, and
the native global-section/spectrum adjunction recovers the literal
original structure morphism. These identities give exactly the original
coefficient triangle for the constructed whole chart isomorphism.
This proves a whole original affine chart comparison; global chart
gluing and the intrinsic model-psh numerical comparison remain separate
producers.
\end{proof}
